\section{Introduction}
\label{sec:intro}

Anonymous credentials (AC)~\cite{chaum85-anonymous-credentials,EC:Brands95,EC:CamLys01} allow parties to
convey certified information about themselves in a privacy-preserving way.  The usual roles are issuers,
holders, and verifiers.  A holder obtains attributes certified by an issuer and later proves statements about
those attributes to a verifier without revealing the credential itself.  The central security goals are
unforgeability and unlinkability: verifiers should only accept credentials derived from an issuer, while issuers
and verifiers should not be able to link a presentation to an issuance session or to another presentation.

Certain AC use cases, such as anonymous rate limiting~\cite{ietf-privacypass-arc-protocol-01,schlesinger-cfrg-act-00},
require that a credential be usable only a bounded number of times in a given context, for example per
epoch or per service.  Existing classical constructions can provide compact credentials and presentations, but
practically efficient post-quantum ACs remain much less developed.  This gap is relevant for digital identity
systems, including wallet deployments that will eventually require post-quantum migration.

\paragraph*{Building ACs.}
Anonymous credentials are often built from a proof-friendly signature and a compatible zero-knowledge proof.
The holder obtains a signature on hidden attributes and later proves, in zero knowledge, knowledge of a valid
message-signature pair.  A generic post-quantum combination of standard signatures, commitments, and
zero-knowledge proofs is currently too large for the settings targeted here.  We therefore combine a proof-friendly
lattice signature, a compact lattice commitment term, and SmallWood zero-knowledge proofs in a relation designed
to remain low degree.

\subsection{Contributions}
We construct a post-quantum issuer-view-hiding issuance design and a conditional anonymous rate-limited
credential design from lattice assumptions and hash-based zero knowledge.  The construction uses proof-friendly
lattice signatures from vanishing SIS with rational hashing~\cite{EPRINT:DKLW25}, an NTRU-style Gaussian
preimage-sampling candidate based on Antrag~\cite{EPRINT:ENSTW23}, and the SmallWood proof system for
relatively small statements~\cite{EPRINT:FenRiv25b}.  The DKLW import is used for the non-blind
hash-and-preimage signature layer.  The interactive issuance protocol is treated separately as a correct and
issuer-view-hiding committed-message issuance layer, with one-more unforgeability left as an explicit assumption
in the ARC soundness decomposition.  Rate limiting is implemented by signing a hidden PRF key and proving correct
tag derivation during presentation.

The issuance layer follows the IntGenISIS/BLNS/DKLW shape.  The holder's semantic signed message is
\[
  M := \vecm\Vert\veck,
\]
where \(\vecm\) is the attribute vector and \(\veck\) is the secret PRF key.  The holder sends a target-shaped
Ajtai/BDLOP-style commitment
\[
  c=C_M M+A_s s+e\in\Rq^{n_c}.
\]
The issuer then samples the BB-tran rational-hash characteristic and randomness independently of \(M\).  With
\(\muSig\in\Rq^{\ell_{\mu_{\mathsf{sig}}}}\), \(x_0\in\Rq^{\ell_{x_0}}\), and \(x_1\in\Rq\), the rational-hash part is
\[
  h_{\muSig,\chiSig}(B)=B_0+B_1\muSig+B_2x_0+Z,
  \qquad
  (B_3-\mathbf 1_n x_1)\Had Z=1_n,
\]
where \(\chiSig=(x_0,x_1)\).  The signed target is
\[
  T=h_{\muSig,\chiSig}(B)+c.
\]
The preimage signature satisfies \(\mA\vecu=T\).  Thus the semantic message appears only in the computationally
hiding commitment term, not directly in the BB-tran rational-hash characteristic.

The issuer-view privacy statement is computational.  It relies on Module-LWE hiding of the commitment term and
zero knowledge of the pre-sign proof; binding of committed openings is reduced to Module-SIS for the block matrix
\([C_M\mid A_s\mid I_{n_c}]\).  This is a local target-shaped specialization of the Ajtai/BDLOP/ABDLOP commitment
precedent used in lattice zero-knowledge protocols~\cite{BDLOP18,LNP22}; we do not claim that this exact compact
form is copied verbatim from those sources.

\subsection{Technical overview}
The semantic signed message is written throughout as
\[
  M=\vecm\Vert\veck,
\]
with an explicit encoding assumption \(\ell_M=1\) for the first target parameter set.  The DKLW rational-hash
message or characteristic is written \(\muSig\) to avoid confusion with the semantic message.

\paragraph*{Issuance.}
The holder samples \(\veck\getsunif\mathcal K_{\mathsf{key}}\), forms \(M=\vecm\Vert\veck\), samples commitment
randomness \(s\leftarrow\chi_s^{k_s}\) and error \(e\leftarrow\chi_e^{n_c}\), and sends
\[
  c=C_M M+A_s s+e
\]
together with a pre-sign NIZK \(\pi_{\mathsf{pre}}\).  This proof shows knowledge of a bounded opening
\((M,s,e)\), correct encoding of \(M=\vecm\Vert\veck\), key-space membership for \(\veck\), and any issuance
policy constraints on \(\vecm\).  It does not prove the BB-tran inverse relation or the final target computation.

The issuer verifies \(\pi_{\mathsf{pre}}\), samples
\[
  \muSig\in\Rq^{\ell_{\mu_{\mathsf{sig}}}},\qquad
  x_0\in\Rq^{\ell_{x_0}},\qquad
  x_1\in\Rq,
\]
rejects and resamples if \(B_3-\mathbf 1_nx_1\) is not invertible, sets
\[
  Z=(B_3-\mathbf 1_nx_1)^{-1},
  \qquad
  T=B_0+B_1\muSig+B_2x_0+Z+c,
\]
and samples \(\vecu\gets\SampPre(\trapdoor,T)\) such that \(\mA\vecu=T\).  The issuer returns
\((\vecu,\muSig,x_0,x_1)\).  The holder stores the credential witness
\[
  (\vecu,M,\vecm,\veck,s,e,\muSig,x_0,x_1).
\]
If \(M\) already carries \(\vecm\) and \(\veck\) in the implementation, \(\vecm\) can be omitted from storage.

\begin{figure}[h]
  \centering
  \resizebox{\columnwidth}{!}{%
  \begin{tikzpicture}[
    actor/.style={font=\bfseries, align=center},
    lifeline/.style={draw, thick},
    msg/.style={-Latex, thick},
    note/.style={draw, rounded corners, fill=gray!10, align=center, inner sep=4pt},
    x=1.15cm, y=0.9cm]
    \node[actor] (I) at (0,10) {\Large Issuer\\[-1pt]$\Issuer$};
    \node[actor] (H) at (7,10) {\Large Holder\\[-1pt]$\Holder$};
    \draw[lifeline] (I) -- ++(0,-8.6);
    \draw[lifeline] (H) -- ++(0,-8.6);
    \node[note] at ($(H)+(0,-1.4)$) {Sample $\veck$ and set $M=\vecm\Vert\veck$\\Sample $s,e$ and compute $c=C_MM+A_s s+e$};
    \draw[msg] ($(H)+(0,-2.7)$) -- node[above,sloped] {$c,\;\pi_{\mathsf{pre}}$} ($(I)+(0,-2.7)$);
    \node[note] at ($(I)+(0,-4.0)$) {Verify $\pi_{\mathsf{pre}}$\\Sample $\muSig,x_0,x_1$\\$Z=(B_3-\mathbf 1_nx_1)^{-1}$};
    \node[note] at ($(I)+(0,-5.8)$) {$T=B_0+B_1\muSig+B_2x_0+Z+c$\\$\vecu\gets\SampPre(\trapdoor,T)$};
    \draw[msg] ($(I)+(0,-7.0)$) -- node[above,sloped] {$\vecu,\muSig,x_0,x_1$} ($(H)+(0,-7.0)$);
    \node[note] at ($(H)+(0,-8.0)$) {Store credential witness\\$(\vecu,M,\vecm,\veck,s,e,\muSig,x_0,x_1)$};
  \end{tikzpicture}}
  \caption{Issuance flow for the committed-message IntGenISIS-style construction.}
  \label{fig:issuance-flow}
\end{figure}

\paragraph*{Showing.}
To present a credential, the holder chooses a nonce, computes
\(\prftag=\PRF(\veck,\nonce)\), reconstructs \(Z=(B_3-\mathbf 1_nx_1)^{-1}\), and proves knowledge of
\(\vecu,M,\vecm,\veck,s,e,\muSig,x_0,x_1,Z\) such that
\begin{align*}
  &(B_3-\mathbf 1_nx_1)\Had Z=1_n,\\
  &\mA\vecu=B_0+B_1\muSig+B_2x_0+Z+C_M M+A_s s+e,\\
  &M=\vecm\Vert\veck,\qquad \prftag=\PRF(\veck,\nonce).
\end{align*}
The verifier sees only \((\nonce,\prftag,\nizkShow)\) and the public parameters.  The issuance commitment \(c\)
is not revealed during showing.

\begin{figure}[h]
  \centering
  \resizebox{\columnwidth}{!}{%
  \begin{tikzpicture}[
    actor/.style={font=\bfseries, align=center},
    lifeline/.style={draw, thick},
    msg/.style={-Latex, thick},
    note/.style={draw, rounded corners, fill=gray!10, align=center, inner sep=4pt},
    x=1.15cm, y=0.9cm]
    \node[actor] (H) at (0,0) {\Large Holder\\[-1pt]$\Holder$};
    \node[actor] (V) at (9,0) {\Large Verifier\\[-1pt]$\Verifier$};
    \draw[lifeline] (H) -- ++(0,-7.1);
    \draw[lifeline] (V) -- ++(0,-7.1);
    \node[note] at ($(H)+(0,-2.4)$) {Compute $\prftag\leftarrow\PRF(\veck,\nonce)$\\Prove knowledge of credential material satisfying\\inverse, signature, commitment-term, encoding, and PRF constraints};
    \draw[msg] ($(H)+(0,-4.3)$) -- node[above,sloped] {$(\nonce,\prftag,\nizkShow)$} ($(V)+(0,-4.3)$);
    \node[note] at ($(V)+(0,-5.8)$) {Reject if $\nizkShow$ fails\\else reject if $(\nonce,\prftag)$ already seen\\else store pair and accept};
    \draw[msg] ($(V)+(0,-6.8)$) -- node[above,sloped] {accept/reject} ($(H)+(0,-6.8)$);
  \end{tikzpicture}}
  \caption{Showing flow.  The proof hides all credential material, including the commitment opening and rational-hash randomness.}
  \label{fig:showing-flow}
\end{figure}

\paragraph*{SmallWood proofs.}
We use SmallWood~\cite{EPRINT:FenRiv25b} as the NIZK for the issuance and showing relations.  The new
relations remain proof-friendly: the commitment equation and the signature equation are linear over the ring,
while the rational-hash inverse equation is bilinear.  The PRF checks retain their existing Poseidon2-style
nonlinear constraints.  Section~\ref{sec:smallwood-model} states the semantic PACS relations and the
coefficient-domain consequences used by the security arguments; Appendix~\ref{app:smallwood} records the
schematic row inventories and proof-stack details.
